\documentclass[12pt,a4paper]{article}
\usepackage[T2A]{fontenc}
\usepackage[utf8]{inputenc}
\usepackage[russian]{babel}
\usepackage[a4paper,margin=2cm]{geometry}
\usepackage{graphicx}
\usepackage{booktabs}
\usepackage{longtable}
\usepackage{array}
\usepackage{amsmath}
\usepackage{amssymb}
\usepackage{listings}
\usepackage{xcolor}
\usepackage{hyperref}
\hypersetup{colorlinks=true,linkcolor=blue,urlcolor=blue}

\lstset{
  basicstyle=\ttfamily\footnotesize,
  breaklines=true,
  frame=single,
  columns=fullflexible,
  keepspaces=true,
  showstringspaces=false,
}

\title{Лабораторная работа \No\,6 \\ Домашнее задание, блок \No\,3 (ДЗ3.1) \\
\large Отбор значимых признаков: градиентный бустинг + UFSACO}
\author{}
\date{}

\begin{document}
\maketitle

\section{Постановка задачи}

Согласно ДЗ3.1, мы воспроизводим пайплайн из \texttt{UFSACO.ipynb} на
собственном наборе данных. Для задачи регрессии нам необходимо:
\begin{enumerate}
  \item с помощью градиентного бустинга (используем как чёрный ящик)
        определить значимые признаки для моделирования целевого;
  \item с помощью метода муравьиной колонии (UFSACO~--- unsupervised
        feature selection) отобрать признаки без использования
        целевого;
  \item обеспечить, чтобы пересечение двух множеств признаков было не
        меньше четырёх элементов; при необходимости варьировать
        гиперпараметры ACO.
\end{enumerate}

Формально: пусть $\mathcal{F} = \{f_1, \dots, f_d\}$~--- полное множество
признаков, $y$~--- целевой признак. Каждый метод возвращает
подмножество
\[
  \mathcal{S}_{\text{GB}} \subseteq \mathcal{F}, \quad
  \mathcal{S}_{\text{ACO}} \subseteq \mathcal{F}, \quad
  |\mathcal{S}_{\text{GB}}| = |\mathcal{S}_{\text{ACO}}| = k,
\]
где $k$~--- число отбираемых признаков; требуется
\[
  |\mathcal{S}_{\text{GB}} \cap \mathcal{S}_{\text{ACO}}| \geq 4,
  \qquad 7 \leq k \leq \lfloor d/2 \rfloor.
\]
Здесь $d$~--- общее число признаков ($d=24$ в нашем случае),
$\lfloor \cdot \rfloor$~--- целая часть.

\section{Описание датасета}

\textbf{Источник:} <<Car Price Prediction Multiple Linear Regression>>,
Kaggle (\href{https://www.kaggle.com/datasets}{kaggle.com/datasets}),
датасет компании Geely Auto, опубликован в~2019~г.

\textbf{Предметная область.} Китайская автомобильная компания Geely
Auto планирует выход на рынок США и хочет понять, какие конструктивные
и потребительские характеристики автомобилей определяют их цену в
американском сегменте. Консалтинговое агентство собрало срез по~205
моделям автомобилей различных производителей. Целевой признак~---
розничная цена в долларах.

\textbf{Размер.} $n = 205$ объектов, $d = 24$ признака~+~целевой
\texttt{price} (всего 25 числовых столбцов после очистки).

\textbf{Целевой признак.} В качестве целевого мы выбрали
$y = \texttt{price}$~--- розничную цену автомобиля в долларах США.
Это вещественный признак, диапазон значений \$5118--45400, среднее
$\bar{y} \approx 13\,277$, стандартное отклонение $\sigma_y \approx 7\,989$.
Выбор обусловлен постановкой исходной задачи на Kaggle (предсказание
цены автомобиля на американском рынке для компании Geely Auto)~---
именно цена является той переменной, факторы влияния на которую мы
хотим понять, что точно соответствует типу задачи в ДЗ
(\emph{регрессия}).

\textbf{Признаки} (после предобработки, см. раздел~\ref{sec:prep}):
\begin{itemize}\setlength\itemsep{0pt}
  \item \emph{Идентификация:} \texttt{carbrand} (марка авто, 22
        уникальных значения, label-encoded).
  \item \emph{Категориальные характеристики кузова и трансмиссии:}
        \texttt{fueltype}, \texttt{aspiration}, \texttt{doornumber},
        \texttt{carbody}, \texttt{drivewheel}, \texttt{enginelocation},
        \texttt{enginetype}, \texttt{fuelsystem}~--- закодированы
        целочисленно.
  \item \emph{Геометрия:} \texttt{wheelbase}, \texttt{carlength},
        \texttt{carwidth}, \texttt{carheight}, \texttt{curbweight}.
  \item \emph{Силовая установка:} \texttt{cylindernumber},
        \texttt{enginesize}, \texttt{boreratio}, \texttt{stroke},
        \texttt{compressionratio}, \texttt{horsepower},
        \texttt{peakrpm}.
  \item \emph{Топливная экономичность:} \texttt{citympg},
        \texttt{highwaympg}.
  \item \emph{Рейтинг страхового риска:} \texttt{symboling}.
  \item \emph{Целевой:} \texttt{price} (диапазон \$5118--45400).
\end{itemize}

Пропусков и дубликатов в данных нет.

\subsection{Предобработка}\label{sec:prep}

Исходный файл \texttt{CarPrice\_Assignment.xls} (фактически CSV
с расширением \texttt{.xls}) приведён к чисто числовой матрице
$X \in \mathbb{R}^{n \times d}$:
\begin{enumerate}\setlength\itemsep{0pt}
  \item удалили служебный \texttt{car\_ID};
  \item из \texttt{CarName} извлекли только марку
        (\texttt{carbrand}); устранили опечатки в названиях марок
        (\texttt{maxda}$\to$\texttt{mazda},
         \texttt{porcshce}$\to$\texttt{porsche},
         \texttt{toyouta}$\to$\texttt{toyota},
         \texttt{vokswagen}/\texttt{vw}$\to$\texttt{volkswagen});
  \item словесные числа в \texttt{doornumber} и
        \texttt{cylindernumber} (\texttt{two}, \texttt{four}, \ldots)
        перевели в целые;
  \item к восьми категориальным признакам применили label-encoding
        (целые от~0 до $k-1$); one-hot не использовали, чтобы остаться
        в рамках $d=24$ признаков и сохранить верхнюю границу отбора
        $\lfloor 24/2 \rfloor = 12$.
\end{enumerate}
Результат сохранили в \texttt{CarPrice\_clean.csv}.

Каждый признак $f_j$ затем приводится к диапазону $[0,1]$ преобразованием
Min--Max:
\begin{equation}
  \tilde{x}_{ij} \;=\; \frac{x_{ij} - \min_i x_{ij}}{\max_i x_{ij} - \min_i x_{ij}},
  \label{eq:minmax}
\end{equation}
где $x_{ij}$~--- значение признака $f_j$ у $i$-го объекта,
$\tilde{x}_{ij}$~--- его масштабированный аналог.

\subsection{Краткий разведочный анализ}

Линейная зависимость каждого признака с целевым оценена через выборочную
корреляцию Пирсона:
\begin{equation}
  r(f_j, y) \;=\; \frac{\sum_{i=1}^{n}\bigl(x_{ij} - \bar{x}_j\bigr)\bigl(y_i - \bar{y}\bigr)}
                       {\sqrt{\sum_{i=1}^{n}\bigl(x_{ij}-\bar{x}_j\bigr)^2}
                        \cdot \sqrt{\sum_{i=1}^{n}\bigl(y_i-\bar{y}\bigr)^2}},
  \label{eq:pearson}
\end{equation}
где $n$~--- число объектов, $x_{ij}$~--- значение $j$-го признака
у $i$-го объекта, $y_i$~--- значение целевого, $\bar{x}_j$ и
$\bar{y}$~--- их средние.

Десять наибольших по модулю корреляций с целевым:

\begin{center}
\begin{tabular}{lr}
\toprule
\textbf{Признак} & \textbf{$r$ с price} \\
\midrule
enginesize        & $+0.874$ \\
curbweight        & $+0.835$ \\
horsepower        & $+0.808$ \\
carwidth          & $+0.759$ \\
cylindernumber    & $+0.718$ \\
highwaympg        & $-0.698$ \\
citympg           & $-0.686$ \\
carlength         & $+0.683$ \\
drivewheel        & $+0.578$ \\
wheelbase         & $+0.578$ \\
\bottomrule
\end{tabular}
\end{center}

Видим, что доминируют <<размерные>> признаки (геометрия~+~масса~+~мощность
двигателя), сильно скоррелированные между собой ($|r| > 0.7$). Этот
факт критически важен~--- он определяет поведение UFSACO ниже.
Расход топлива \texttt{highwaympg}/\texttt{citympg} коррелирует с ценой
отрицательно (чем больше миль на галлон, тем меньше/легче машина, тем
дешевле).

Гистограммы распределений сохранены в \texttt{hist\_hw3.png},
тепловая карта корреляций~--- в \texttt{corr\_heatmap\_hw3.png}.

\section{Метод 1: Gradient Boosting}

\subsection{Модель}
Используется \texttt{sklearn.ensemble.GradientBoostingRegressor}.
Модель строит аддитивную композицию из $M$ регрессионных деревьев $h_m$:
\begin{equation}
  \hat{F}_M(\mathbf{x}) \;=\; F_0 \;+\; \nu \sum_{m=1}^{M} h_m(\mathbf{x}),
  \label{eq:gb}
\end{equation}
где $\mathbf{x}$~--- вектор признаков объекта,
$F_0$~--- начальное предсказание (среднее по $y$),
$\nu \in (0,1]$~--- темп обучения (learning rate),
$M$~--- число деревьев,
$h_m$~--- $m$-е дерево, обучаемое на градиенте функции потерь от
предсказаний предыдущей композиции $\hat{F}_{m-1}$.

В качестве функции потерь используется квадратичная:
\begin{equation}
  L(y, \hat{F}) \;=\; \tfrac{1}{2}\bigl(y - \hat{F}\bigr)^2,
\end{equation}
тогда антиградиент совпадает с остатком $y - \hat{F}$, и на каждой
итерации дерево $h_m$ приближает именно остаток предыдущей модели.

\subsection{Гиперпараметры}
\begin{center}
\begin{tabular}{lr}
\toprule
\texttt{n\_estimators} ($M$)      & 300  \\
\texttt{max\_depth}               & 4    \\
\texttt{min\_samples\_split}      & 10   \\
\texttt{learning\_rate} ($\nu$)   & 0.01 \\
\texttt{random\_state}            & 42   \\
\bottomrule
\end{tabular}
\end{center}

\subsection{Значимость признаков}
Для каждого признака $f_j$ значимость определяется как нормированная
сумма уменьшений функции потерь во всех узлах всех деревьев, где этот
признак использован для расщепления:
\begin{equation}
  I(f_j) \;=\; \frac{1}{M}\sum_{m=1}^{M}\sum_{t \in \mathcal{T}_m^{(j)}}
                 \omega_t \cdot \Delta L_t,
  \label{eq:fimp}
\end{equation}
где $\mathcal{T}_m^{(j)}$~--- множество узлов $m$-го дерева, где
расщепление сделано по признаку $f_j$, $\omega_t = n_t/n$~--- доля
объектов, попавших в узел $t$, $\Delta L_t$~--- уменьшение функции
потерь в этом узле, $n_t$~--- размер выборки узла, $n$~--- общий размер
обучающей выборки.

После обучения качество на обучающей выборке:
\[
  R^2 \;=\; 1 - \frac{\sum_{i=1}^{n}(y_i - \hat{y}_i)^2}{\sum_{i=1}^{n}(y_i - \bar{y})^2}
        \;=\; 0.9813,
\]
где $\hat{y}_i$~--- предсказание модели для $i$-го объекта.

Сортируем $I(f_j)$ по убыванию и берём top-$k$ при $k=12$ (см.
раздел~\ref{sec:n}):

\begin{center}
\begin{tabular}{cl}
\toprule
\textbf{Ранг} & \textbf{Признак} \\
\midrule
1  & enginesize \\
2  & curbweight \\
3  & highwaympg \\
4  & horsepower \\
5  & carlength \\
6  & carwidth \\
7  & carbrand \\
8  & wheelbase \\
9  & stroke \\
10 & compressionratio \\
11 & peakrpm \\
12 & fuelsystem \\
\bottomrule
\end{tabular}
\end{center}

См. также \texttt{gb\_importance\_hw3.png}.

\section{Метод 2: UFSACO (муравьиная колония)}

\subsection{Идея}
Используется реализация из \texttt{UFSACO.ipynb} без изменений. Алгоритм
\emph{unsupervised}: целевой признак не участвует в построении графа.

Признаки рассматриваются как вершины полного неориентированного графа
$G = (\mathcal{F}, E)$. Каждому ребру $(i,j)$ присваивается вес, равный
\emph{косинусной близости} векторов признаков:
\begin{equation}
  s(f_i, f_j) \;=\; \frac{\langle \mathbf{x}^{(i)}, \mathbf{x}^{(j)} \rangle}
                          {\|\mathbf{x}^{(i)}\| \cdot \|\mathbf{x}^{(j)}\|},
  \label{eq:cos}
\end{equation}
где $\mathbf{x}^{(i)} \in \mathbb{R}^n$~--- столбец масштабированной
матрицы признаков, соответствующий признаку $f_i$,
$\langle \cdot, \cdot \rangle$~--- скалярное произведение, $\|\cdot\|$~---
евклидова норма. Высокое $s(f_i,f_j)$ означает высокую избыточность
этой пары.

Каждой вершине $j$ соответствует уровень феромона $\tau_j \geq 0$,
изначально $\tau_j = \tau_0$.

\subsection{Шаг муравья}
Если $k$-й муравей находится в вершине $i$ и $V_k$~--- множество уже
посещённых им вершин, то <<желательность>> следующего перехода в
вершину $j \in \mathcal{F} \setminus (V_k \cup \{i\})$:
\begin{equation}
  \eta_{ij} \;=\; \frac{\tau_j}{\bigl(s(f_i, f_j) + \varepsilon\bigr)^{\alpha}},
  \label{eq:eta}
\end{equation}
где $\tau_j$~--- текущий уровень феромона на вершине $j$,
$s(f_i, f_j)$~--- косинусная близость (\ref{eq:cos}),
$\alpha > 0$~--- параметр относительной важности
<<некоррелированности>>,
$\varepsilon$~--- малая константа для численной устойчивости.
Чем меньше $s(f_i,f_j)$ и больше $\tau_j$, тем выгоднее переход.

С вероятностью $q_0 \in [0,1]$ (\texttt{EXPLOITATION\_PROB}) муравей
переходит детерминированно (exploitation):
\begin{equation}
  j^* \;=\; \arg\max_{j \notin V_k \cup \{i\}} \eta_{ij},
\end{equation}
иначе (exploration)~--- стохастически с вероятностями,
пропорциональными $\eta_{ij}$:
\begin{equation}
  P(j \mid i, V_k) \;=\;
    \frac{\eta_{ij}}{\sum_{\ell \notin V_k \cup \{i\}} \eta_{i\ell}}.
  \label{eq:explore}
\end{equation}

\subsection{Обновление феромона}
После того как все $N_{\text{ants}}$ муравьёв сделали по
$N_{\text{steps}}$ шагов внутри одной эпохи, для каждой вершины
$j$ известно количество визитов $v_j$. Феромон обновляется по правилу:
\begin{equation}
  \tau_j \;\leftarrow\; (1 - \rho)\,\tau_j \;+\; \frac{v_j}{\sum_{\ell} v_\ell},
  \label{eq:tau}
\end{equation}
где $\rho \in (0,1]$~--- коэффициент испарения (\texttt{RO}),
$v_j$~--- число визитов вершины $j$ за текущую эпоху,
$\sum_\ell v_\ell$~--- общее число визитов (нормирует прирост).

\subsection{Финальный отбор}
Внешний цикл повторяется $\text{NC}_{\max}$ эпох. По окончании
ранжируем вершины по $\tau_j$ и берём $k$ с наибольшими значениями:
\begin{equation}
  \mathcal{S}_{\text{ACO}} \;=\; \mathrm{Top}_k\!\bigl\{\tau_1, \dots, \tau_d\bigr\}.
\end{equation}

\subsection{Замечание о масштабировании}
В ноутбуке применён \texttt{MinMaxScaler}~(\ref{eq:minmax}), но в нашем
датасете несколько закодированных категориальных признаков
(\texttt{enginelocation}, \texttt{aspiration}, \texttt{doornumber})
почти константы~--- одно значение встречается у $\geq 90\%$ объектов.
После Min--Max их вектор $\mathbf{x}^{(j)}$ близок к нулевому,
норма $\|\mathbf{x}^{(j)}\| \to 0$, и косинусная близость
$s(f_i,f_j)$ становится численно неустойчивой~--- такие признаки
получают завышенные $\eta_{ij}$ и захватывают весь top-$k$.

Чтобы устранить эту патологию, мы дополнительно центрируем матрицу
признаков преобразованием Z-score:
\begin{equation}
  z_{ij} \;=\; \frac{x_{ij} - \bar{x}_j}{\sigma_j},
  \label{eq:zscore}
\end{equation}
где $\bar{x}_j$ и $\sigma_j$~--- выборочное среднее и стандартное
отклонение $j$-го признака. На центрированных векторах формула
(\ref{eq:cos}) переходит в коэффициент корреляции Пирсона~--- ту же
меру~(\ref{eq:pearson}), что используется для EDA. Сам алгоритм UFSACO
при этом остаётся без изменений.

\subsection{Гиперпараметры}\label{sec:hp}
Дефолтные значения из ноутбука дали слишком малое пересечение с GB
(см. раздел~\ref{sec:tune}). Финальные значения:

\begin{center}
\begin{tabular}{ll r}
\toprule
Обозначение & Параметр & Значение \\
\midrule
$\text{NC}_{\max}$        & число эпох (\texttt{NC\_MAX})            & 50   \\
$N_{\text{steps}}$        & шагов в эпохе (\texttt{N\_STEPS})        & 10   \\
$N_{\text{ants}}$         & число муравьёв (\texttt{N\_ANTS})        & 25   \\
$\tau_0$                  & начальный феромон (\texttt{INIT\_PHER.}) & 0.2  \\
$\rho$                    & коэф.\,испарения (\texttt{RO})           & 0.1  \\
$q_0$                     & вероятность exploitation                 & 0.7  \\
$\alpha$                  & вес <<некоррелированности>>              & 0.2  \\
$k$                       & число отбираемых признаков               & 12   \\
\bottomrule
\end{tabular}
\end{center}

\subsection{Устойчивость: 5 запусков}

Чтобы оценить разброс, мы запустили алгоритм 5 раз; финальный
список~--- 12 признаков, чаще всего попадавших в top-12 (по частоте
появлений). Если обозначить $c_j$~--- число запусков, в которых
признак $f_j$ попал в top-$k$, то:
\begin{equation}
  \mathcal{S}_{\text{ACO}}^{\text{final}}
    \;=\; \mathrm{Top}_k\bigl\{c_1, \dots, c_d\bigr\}.
\end{equation}

Итоговый top-12 UFSACO:
\texttt{enginetype, stroke, cylindernumber, doornumber, carbody,
peakrpm, symboling, boreratio, aspiration, highwaympg, horsepower,
compressionratio}.

\section{Сравнение и пересечение}

\begin{longtable}{lcc}
\toprule
\textbf{Признак} & \textbf{GB} & \textbf{UFSACO} \\
\midrule
\endhead
aspiration        &              & $\checkmark$ \\
boreratio         &              & $\checkmark$ \\
carbody           &              & $\checkmark$ \\
carbrand          & $\checkmark$ &              \\
carlength         & $\checkmark$ &              \\
carwidth          & $\checkmark$ &              \\
\textbf{compressionratio} & $\checkmark$ & $\checkmark$ \\
curbweight        & $\checkmark$ &              \\
cylindernumber    &              & $\checkmark$ \\
doornumber        &              & $\checkmark$ \\
enginesize        & $\checkmark$ &              \\
enginetype        &              & $\checkmark$ \\
fuelsystem        & $\checkmark$ &              \\
\textbf{highwaympg} & $\checkmark$ & $\checkmark$ \\
\textbf{horsepower} & $\checkmark$ & $\checkmark$ \\
\textbf{peakrpm}    & $\checkmark$ & $\checkmark$ \\
\textbf{stroke}     & $\checkmark$ & $\checkmark$ \\
symboling         &              & $\checkmark$ \\
wheelbase         & $\checkmark$ &              \\
\bottomrule
\end{longtable}

\textbf{Пересечение:}
\[
  \mathcal{S}_{\text{GB}} \cap \mathcal{S}_{\text{ACO}}
    = \{\texttt{compressionratio}, \texttt{highwaympg}, \texttt{horsepower},
        \texttt{peakrpm}, \texttt{stroke}\}.
\]
Мощность пересечения
$|\mathcal{S}_{\text{GB}} \cap \mathcal{S}_{\text{ACO}}| = 5 \geq 4$~---
требование выполнено.

Дополнительно мы оценили <<согласованность>> двух методов через
коэффициент Жаккара:
\begin{equation}
  J(\mathcal{S}_{\text{GB}}, \mathcal{S}_{\text{ACO}}) \;=\;
    \frac{|\mathcal{S}_{\text{GB}} \cap \mathcal{S}_{\text{ACO}}|}
         {|\mathcal{S}_{\text{GB}} \cup \mathcal{S}_{\text{ACO}}|}
    \;=\; \frac{5}{19} \;\approx\; 0.263,
\end{equation}
где $|\cdot|$~--- мощность множества. Значение $J \in [0,1]$;
$J = 1$ означает полное совпадение, $J = 0$~--- непересечение.

\subsection{Выбор $k$}\label{sec:n}
По правилу ЛР, $7 \leq k \leq \lfloor 24/2 \rfloor = 12$.
Эксперименты при $k = 7,\,10,\,12$:

\begin{center}
\begin{tabular}{ccl}
\toprule
$k$ & $|\mathcal{S}_{\text{GB}} \cap \mathcal{S}_{\text{ACO}}|$ & комментарий \\
\midrule
7   & 0--2 & пересечение мало даже при тюнинге \\
10  & 2    & лучше, но не до~4 \\
\textbf{12} & \textbf{5} & требование выполнено \\
\bottomrule
\end{tabular}
\end{center}

\subsection{Подбор гиперпараметров ACO}\label{sec:tune}
С дефолтами ноутбука ($\text{NC}_{\max}=3$, $N_{\text{steps}}=4$,
$N_{\text{ants}}=5$, $\alpha=1$, \texttt{MinMaxScaler}) при $k=7$
пересечение составило~2. Причина в характере данных: GB-top~--- единый
кластер сильно скоррелированных <<размерных>> признаков
(\texttt{enginesize}, \texttt{curbweight}, \texttt{carlength},
\texttt{carwidth}, \texttt{horsepower}), а UFSACO по дизайну, согласно
формуле~(\ref{eq:eta}), штрафует за высокое $s(f_i,f_j)$~--- то есть
выбирает по одному представителю каждого кластера. Чтобы поднять
перекрытие, мы:
\begin{enumerate}\setlength\itemsep{0pt}
  \item увеличили $k$ до~12 (тогда GB захватывает не только <<размерный>>
        кластер, а UFSACO~--- больше представителей других кластеров);
  \item понизили $\alpha$ с~1.0 до~0.2; согласно~(\ref{eq:eta}), при
        $\alpha \to 0$ знаменатель $\to 1$ и штраф за сходство
        ослабевает, что позволяет ACO выбирать несколько членов одного
        кластера;
  \item увеличили число эпох ($\text{NC}_{\max}$: 3$\to$50), шагов
        ($N_{\text{steps}}$: 4$\to$10) и муравьёв ($N_{\text{ants}}$:
        5$\to$25) для стабилизации феромона $\tau$ в (\ref{eq:tau});
  \item понизили испарение $\rho$ с~0.2 до~0.1, чтобы накопленный опыт
        переходов не <<вымывался>> слишком быстро (множитель $1-\rho$ в
        (\ref{eq:tau})).
\end{enumerate}

\section{Интерпретация результатов}

Пять общих признаков логично интерпретируются:
\begin{itemize}\setlength\itemsep{0pt}
  \item \texttt{horsepower}~--- мощность; ключевой драйвер цены, оба
        метода согласны.
  \item \texttt{highwaympg}~--- экономичность; обратно связана с
        мощностью, но не полностью~--- UFSACO видит в ней
        самостоятельный кластер.
  \item \texttt{peakrpm}~--- характеристика оборотистости двигателя,
        слабо коррелирует с размерным кластером.
  \item \texttt{stroke}~--- ход поршня, конструктивная характеристика
        двигателя, малая корреляция со всем остальным.
  \item \texttt{compressionratio}~--- степень сжатия, аналогично~---
        отдельный кластер.
\end{itemize}

Признаки <<размерного>> кластера (\texttt{enginesize},
\texttt{curbweight}, \texttt{carlength}, \texttt{carwidth},
\texttt{carbrand}, \texttt{wheelbase}) попали только в GB-top:
ACO выбирает максимум одного представителя такого кластера и переключается
на менее коррелированные признаки. Это нормальное и предсказуемое
поведение unsupervised-метода~--- его задача согласно~(\ref{eq:eta})
найти \emph{разнообразное} подпространство, а не максимально
предсказательное.

\section{Воспроизводимость}

\textbf{Код:} \texttt{ufsaco\_hw3.py} (один файл, запуск
\texttt{python ufsaco\_hw3.py}).

\textbf{Зависимости:} \texttt{numpy, pandas, matplotlib, seaborn,
scikit-learn}.

\textbf{Артефакты после запуска:}
\begin{itemize}\setlength\itemsep{0pt}
  \item \texttt{CarPrice\_clean.csv}~--- предобработанный датасет;
  \item \texttt{eda\_describe\_hw3.csv}~--- описательные статистики;
  \item \texttt{hist\_hw3.png}~--- гистограммы;
  \item \texttt{corr\_with\_price\_hw3.csv},
        \texttt{corr\_heatmap\_hw3.png}~--- корреляции;
  \item \texttt{gb\_importance\_hw3.png}~--- важности GB;
  \item \texttt{sim\_top10\_hw3.csv}, \texttt{sim\_bot10\_hw3.csv}~---
        самые похожие/различные пары признаков по cos-similarity;
  \item \texttt{results\_hw3.json}~--- финальные результаты в
        машиночитаемом виде.
\end{itemize}

\section{Выводы}

\begin{enumerate}\setlength\itemsep{0pt}
  \item Датасет CarPrice (205~$\times$~24) полностью удовлетворяет
        требованиям ДЗ к составу, размеру и году публикации.
  \item Градиентный бустинг согласно~(\ref{eq:gb}) достигает
        $R^2 = 0.98$ на обучении и ожидаемо выделяет <<размерный>>
        кластер как наиболее информативный для прогноза цены.
  \item UFSACO ведёт себя как unsupervised-метод: согласно
        формуле~(\ref{eq:eta}) выбирает \emph{взаимно различные}
        признаки, а не самые предсказательные. На датасете с одним
        доминирующим скоррелированным кластером это естественно
        расходится с GB.
  \item После настройки гиперпараметров и выбора $k=12$ мы получили
        пересечение из 5~признаков
        (\texttt{compressionratio, highwaympg, horsepower, peakrpm,
        stroke}), что превосходит порог в~4 признака; коэффициент
        Жаккара $J \approx 0.263$.
\end{enumerate}

\appendix
\section{Top-12 GB по убыванию важности}

\begin{center}
\includegraphics[width=0.75\linewidth]{gb_importance_hw3.png}
\end{center}

\section{Корреляционная матрица}

\begin{center}
\includegraphics[width=0.95\linewidth]{corr_heatmap_hw3.png}
\end{center}

\end{document}
